\documentclass[10pt,a4paper]{amsart}
\usepackage[margin=19mm]{geometry}
\usepackage{amsmath,amssymb,mathtools}
\usepackage[scr=rsfs,cal=euler]{mathalfa}
\usepackage{xfrac}
\usepackage[unicode,hidelinks,bookmarks=false]{hyperref}
\newcommand{\ie}{i.e.\ }
\newcommand{\cf}{cf.\ }
\newcommand{\resp}{respectively\ }
\newcommand{\Subsection}{\subsection}
\providecommand{\nobreakdash}{\nobreak\hbox{-}}
\setlength{\emergencystretch}{2em}
\multlinegap0pt
\usepackage{bm}
\usepackage{bbm}
\usepackage{dsfont}
\usepackage{xspace}
\usepackage{cancel}
\usepackage{mathtools}
\usepackage{amsthm}
\makeatletter
\@ifundefined{theorem}{\newtheorem{theorem}{Theorem}}{}
\makeatother
\makeatletter
\newcommand{\argmin}{\mathop{\mathrm{arg min}}}
\newcommand{\calF}{\mathcal{F}}
\newcommand{\calH}{\mathcal{H}}
\newcommand{\dd}{\mathrm{d}}
\expandafter\ifx\csname proof\endcsname\relax
\renewenvironment{proof}{\par\noindent{\bf Proof\ }}{\hfill\BlackBoxProof\par\medskip}
\fi
\newcommand{\prox}{\mathrm{prox}}
\makeatother
\title[Global Convergence and Error Propagation in Neural Gradient ]{Global Convergence and Error Propagation in Neural Gradient Flows: A Riemannian Optimization Framework}
\author{Shixin Zheng, Yiwei Wang, Haizhao Yang}
\date{Source: arXiv:2605.27779v1 (CC BY 4.0)}
\begin{document}
\maketitle

\noindent\emph{Source and licence.} Global Convergence and Error Propagation in Neural Gradient Flows: A Riemannian Optimization Framework. Version arXiv:2605.27779v1 (CC BY 4.0), distributed under the Creative Commons Attribution 4.0 International licence, as recorded in the arXiv licence metadata for this exact version and shown on the abstract page https://arxiv.org/abs/2605.27779v1. Full rights notice: licenses/arxiv-2605.27779-CC-BY-4.0.txt.\par\smallskip

\noindent\textit{Editorial scope.} Counted unit: the Theorem whose source label is \texttt{thm:mms-convergence} in the article (source0.tex lines 3290--3319, source label \texttt{thm:mms-convergence}), delivered together with the article's own complete original proof of it (source0.tex lines 3322--3468, one \texttt{proof} environment of 147 lines). No mathematics is written, repaired or re-derived: statement and proof are the article's own text line for line. Required context retained uncounted: none. Every definition, display and label of those lines is retained unchanged. Omitted from this package and not redistributed: 1--3289, 3320--3321, 3469--4333. Nothing inside a retained excerpt is omitted or altered: every retained body is delivered exactly as the article's lines are written, so no omission receipt of any kind accompanies this package. Citations: the retained excerpts carry no citation command at all, so no bibliography entry of the article is transcribed and no reference is redistributed. Reference adaptation: none was needed and none was made; every reference inside the delivered unit resolves inside it, so no unresolved reference or citation is produced. Printed slips kept literal: no printed word, symbol, hypothesis or number of the counted statement or of its proof was altered. Layout and class: the article's own class is \texttt{article} with packages including amsmath, amsthm, jmlr2e, subcaption, enumitem, mathtools, bm, overpic, booktabs, makecell, lineno, multirow, ulem, xcolor; the wrapper is the lane's pinned \texttt{amsart} editorial wrapper at 10pt on A4, which restates the theorem-like environments the retained text needs and adds the article's own macro definitions that the retained text needs (\texttt{\textbackslash argmin}, \texttt{\textbackslash calF}, \texttt{\textbackslash calH}, \texttt{\textbackslash dd}, \texttt{\textbackslash prox}), with extra wrapper packages (bm, bbm, dsfont, xspace, cancel, mathtools). The delivered numbering is the unit's own. Assets: no retained span contains a figure environment, a graphics-inclusion command, a PDF-page inclusion, a table or a TikZ picture; no image file is copied into this package, the package declares no asset dependency and no image file is redistributed with it. Licence: version arXiv:2605.27779v1 of this article is distributed under the Creative Commons Attribution 4.0 International licence, as recorded in the arXiv licence metadata for this exact version and shown on the abstract page https://arxiv.org/abs/2605.27779v1. A full rights notice is delivered at licenses/arxiv-2605.27779-CC-BY-4.0.txt. Characters: every retained excerpt is pure ASCII, so the delivered engine, XeLaTeX with its default Latin Modern text font, reports no missing character.
\begin{theorem}\label{thm:mms-convergence}
Let $\calH$ be a Hilbert space, and let $\calF:\calH\to\mathbb{R}$ be continuously Fr\'echet differentiable, coercive, and $m_{\calF}$-strongly convex with $m_\calF>0$. Assume moreover that, for every $E\in\mathbb{R}$, the sublevel set
\[
S_E^{\calF}:=\{u\in\calH:\calF[u]\le E\}
\]
is compact in $\calH$. Given $u^0\in \mathrm{dom}(\calF)$ and $\tau>0$, define the minimizing movement scheme by
\begin{equation}\label{eq:appendix-mms}
u^{n+1}=J_\tau(u^n),
\qquad
J_\tau(x):=\prox_{\tau\calF}(x)
:=\argmin_{u\in\calH}
\left\{
\frac{1}{2\tau}\|u-x\|_{\calH}^2+\calF[u]
\right\}.
\end{equation}
If $1+\tau m_{\calF}>0$, then the minimization problem in \eqref{eq:appendix-mms} admits a unique solution, making $J_\tau$ well defined and single-valued.

Let $\hat u_\tau:[0,T]\to\calH$ denote the piecewise linear interpolation associated with the sequence $\{u^n\}_{n\ge0}$. As $\tau\downarrow0$, the interpolants $\hat u_\tau$ converge in $\calH$, uniformly on $[0,T]$, to the unique solution
\(
u\in W^{1,2}([0,T];\calH) 
\)
of the gradient flow equation $u_t = - \nabla \calF[u]$ associated with $\calF$ and the initial datum $u^0$. Moreover, for any two gradient flow solutions $u_0,u_1$, one has the $m_{\calF}$-contractivity estimate
\begin{equation*}
\|u_1(t)-u_0(t)\|_{\calH}
\le
e^{-m_{\calF}(t-s)}
\|u_1(s)-u_0(s)\|_{\calH},
\qquad 0\le s<t.
\end{equation*}
\end{theorem}

\begin{proof} 

\noindent \textbf{Step 1: Well-posedness and Energy Estimate.} The objective function in \eqref{eq:appendix-mms} is given by $\Phi(u) = \frac{1}{2\tau}\|u-u^n\|_{\calH}^2+\calF[u]$. Since $\calF$ is continuously Fr\'echet differentiable, coercive, and has compact sublevel sets, $\Phi(u)$ is coercive and bounded from below, guaranteeing the existence of a minimizer. Furthermore, because $\calF$ is $m_{\calF}$-convex, $\Phi(u)$ is strictly convex provided $1 + \tau m_{\calF} > 0$. Thus, $u^{n+1}$ exists and is unique. Moreover, $\Phi(u^{n+1}) \le \Phi(u^n)$ indicates
\begin{equation}
 \frac{1}{2 \tau} \| u^{n+1} - u^n \|^2 \leq \mathcal{F}[u^{n}] - \mathcal{F}[u^{n+1}]
\end{equation}
Summing this inequality over $n=0, \dots, N-1$, we have
\begin{equation}\label{sum_incre}
 \sum_{n=0}^{N-1}\| u^{n+1} - u^n \|^2 \leq 2 \tau (\mathcal{F}(u^0) - \mathcal{F}(u^*) ) \leq \tau C  ,
\end{equation}
where $u^*$ is the global minimizer of $\mathcal{F}$. Hence,
\begin{equation}
  \| u^{n+1} - u^n \|^2 \rightarrow 0, \quad n \rightarrow \infty
\end{equation}


\noindent \textbf{Step 2: Interpolations and Compactness.} Define the piecewise affine interpolation $\hat u_\tau:[0,T]\to\mathcal H$ and the piecewise constant interpolation $\bar u_\tau:[0,T]\to\mathcal H$ by
\[
\hat u_\tau((n+\theta)\tau)
=(1-\theta)u^n+\theta u^{n+1},
\qquad \theta\in[0,1],\quad n=0,\dots,N-1,
\]
and
\[
\bar u_\tau(t)=u^{n+1},
\qquad t\in(n\tau,(n+1)\tau],\quad n=0,\dots,N-1.
\]
Then, (\ref{sum_incre}) indicates that
\begin{equation}
\int_{0}^T  \| \dot{\hat{u}}_{\tau} \|^2 \dd t = \sum_{n=1}^N \tau \left \| \frac{u^{n+1} - u^n}{\tau} \right \|^2 \leq 2 (\calF[u^0] - \calF[u^*])
\end{equation}
Hence, $\hat{u}_{\tau}$ is bounded in $W^{1, 2} ([0, T], \calH)$ and we can extract a subsequence (not relabeled) such that
\begin{equation}
\hat{u}_{\tau} \rightharpoonup u ~\text{in}~ L^2([0, T]; \calH), \quad \dot{\hat{u}}_{\tau} \rightharpoonup \dot{u} ~\text{in}~ L^2([0, T]; \calH)
\end{equation}
for a limit $u \in W^{1, 2} ([0, T]; \calH)$.

Moreover, for any $0\le s<t\le T$, since $\hat u_\tau$ is absolutely continuous on $[0,T]$, we have
\(
\hat u_\tau(t)-\hat u_\tau(s)=\int_s^t \dot{\hat u}_\tau(r)\,dr.
\)
Therefore, by the Cauchy--Schwarz inequality, we have
\[
\|\hat u_\tau(t)-\hat u_\tau(s)\|_{\mathcal H}
\le \int_s^t \|\dot{\hat u}_\tau(r)\|_{\mathcal H}\,dr
\le |t-s|^{1/2}\|\dot{\hat u}_\tau\|_{L^2([0,T];\mathcal H)} \le |t-s|^{1/2} \bigl(2(\mathcal F[u^0]-\calF[u^*])\bigr)^{1/2}
\]
Hence $\{\hat u_\tau\}_\tau$ is uniformly H\"older continuous on every $[0,T]$. Since $\hat{u}_{\tau}$ lie in the compact sublevel set $S_{\mathcal{F}(u^0)}^\calF$, by the Arzel\`a--Ascoli theorem, by extracting a further subsequence (not relabeled), we have the uniform convergence
\begin{equation}
 \hat{u}_{\tau} \rightarrow u ~\text{in}~ C^0([0, T]; \mathcal{H}) 
\end{equation}
Notice $\hat{u} (n \tau) = \bar{u} (n \tau)$, and $\bar{u}_{\tau}$ is piecewise constant, we have
\begin{equation}
 \| \hat{u}_{\tau} - \bar{u}_{\tau} \|_{L^{\infty} ([0, T]; \calH)} \leq \sqrt{\tau} \bigl(2(\mathcal F[u^0]-\calF[u^*])\bigr)^{1/2} \rightarrow 0, \quad \tau \rightarrow 0^+.
\end{equation}
Hence, for the piecewise constant interpolation, we have
\begin{equation}
  \bar{u}_{\tau} \rightarrow u ~\text{in}~ L^{\infty} ([0, T]; \mathcal{H})
\end{equation}

\noindent \textbf{Step 3: Convergence to the Gradient Flow.} Next, we show that the limit $u$ satisfies the gradient flow equation
\begin{equation}
u_t = - \nabla \mathcal{F}(u)
\end{equation}
Since $\mathcal F$ is Fr\'echet continuously differentiable, the Euler--Lagrange equation for the minimizing movement scheme reads
\[
\frac{u^{n+1}-u^n}{\tau}=-\nabla\mathcal F[u^{n+1}],
\qquad n=0,1,\dots,N-1.
\]
Recalling the piecewise affine interpolation $\hat u_\tau$ and the piecewise constant interpolation $\bar u_\tau$, we have
\begin{equation}\label{eq:gf-discrete-interpolant}
\dot{\hat u}_\tau(t)=-\nabla\mathcal F[\bar u_\tau(t)]
\qquad \text{for a.e. } t\in(0,T).
\end{equation}
We have proved that up to a subsequence,
\[
\hat u_\tau \to u \qquad \text{in } C^0([0,T];\mathcal H), \dot{\hat u}_\tau \rightharpoonup \dot u
\qquad \text{weakly in } L^2([0,T];\mathcal H), \bar u_\tau \to u
\qquad \text{in } L^\infty([0,T];\mathcal H).
\]
Since the trajectories $\bar u_\tau(t)$ remain in the compact sublevel set
\(
S_{\mathcal F[u^0]}^{\mathcal F} \)
and since $\nabla\mathcal F$ is continuous on bounded subsets of $\mathcal H$, it follows that
\[
\nabla\mathcal F[\bar u_\tau]\to \nabla\mathcal F[u]
\qquad \text{in } L^\infty([0,T];\mathcal H),
\]
and therefore also in $L^2([0,T];\mathcal H)$. We now pass to the limit in \eqref{eq:gf-discrete-interpolant}. For any $\phi\in L^2([0,T];\mathcal H)$,
\[
\int_0^T \langle \dot{\hat u}_\tau(t),\phi(t)\rangle_{\mathcal H}\,dt
=
-\int_0^T \langle \nabla\mathcal F[\bar u_\tau(t)],\phi(t)\rangle_{\mathcal H}\,dt.
\]
Letting $\tau\to0$, the weak convergence of $\dot{\hat u}_\tau$ and the strong convergence of $\nabla\mathcal F[\bar u_\tau]$ yield
\[
\int_0^T \langle \dot u(t),\phi(t)\rangle_{\mathcal H}\,dt
=
-\int_0^T \langle \nabla\mathcal F[u(t)],\phi(t)\rangle_{\mathcal H}\,dt.
\]
Since $\phi$ is arbitrary, we conclude that
\[
\dot u=-\nabla\mathcal F[u]
\qquad \text{in } L^2([0,T];\mathcal H).
\]
That is, $u$ satisfies the gradient flow equation on $[0,T]$. Finally, since $\hat u_\tau(0)=u^0$ for all $\tau$ and $\hat u_\tau\to u$ in $C([0,T];\mathcal H)$, we also have
\(
u(0)=u^0.
\)
Therefore $u$ is a solution of the gradient flow equation with initial datum $u^0$.


\noindent \textbf{Step 4: Contractivity and Uniqueness.} Since $\mathcal F:\mathcal H\to\mathbb R$ is $m_{\mathcal F}$-convex, i.e.,
\[
\mathcal F[v]\ge \mathcal F[u]
+\langle \nabla\mathcal F[u],\,v-u\rangle_{\mathcal H}
+\frac{m_{\mathcal F}}{2}\|v-u\|_{\mathcal H}^2,
\qquad \forall u,v\in\mathcal H.
\]
Then the gradient $\nabla\mathcal F$ is $m_{\mathcal F}$-monotone:
\begin{equation}\label{eq:monotone-gradient}
\langle \nabla\mathcal F[u]-\nabla\mathcal F[v],\,u-v\rangle_{\mathcal H}
\ge m_{\mathcal F}\|u-v\|_{\mathcal H}^2,
\qquad \forall u,v\in\mathcal H.
\end{equation}
Let $u_0,u_1\in W^{1,2}([0,T];\mathcal H)$ be two solutions of the gradient flow equation
\[
\dot u_j(t)=-\nabla\mathcal F[u_j(t)],
\qquad j=0,1.
\]
Since $u_0,u_1\in W^{1,2}([0,T];\mathcal H)$, using the gradient flow equation and \eqref{eq:monotone-gradient}, we obtain
\[
\begin{aligned}
\frac12\frac{d}{dt}\|u_1(t)-u_0(t)\|_{\mathcal H}^2
&=
-\langle \nabla\mathcal F[u_1(t)]-\nabla\mathcal F[u_0(t)],\,u_1(t)-u_0(t)\rangle_{\mathcal H} \\
&\le
-\,m_{\mathcal F}\|u_1(t)-u_0(t)\|_{\mathcal H}^2 .
\end{aligned}
\]
By Gr\"onwall's inequality, for all $0\le s<t\le T$,
\begin{equation*}
\|u_1(t)-u_0(t)\|_{\mathcal H}
\le e^{-m_{\mathcal F}(t-s)}\|u_1(s)-u_0(s)\|_{\mathcal H}.
\end{equation*}
In particular, if $u_1(0)=u_0(0)=u^0$, then $u_1\equiv u_0$ on $[0,T]$. Hence the gradient flow solution with initial datum $u^0$ is unique.
\end{proof}

\end{document}
